\section{Block Cipher, DES, dan AES}

\subsection{Block Cipher}
Block cipher mengenkripsi data dalam blok-blok tetap (misalnya 64 atau 128 bit). Algoritma seperti DES dan AES adalah block cipher \cite{stallings2017,schneier1996}. Operasi dasar meliputi substitusi (S-box), permutasi (P-box), dan operasi XOR dengan subkunci.

\begin{table}[!htbp]
\centering
\caption{Perbandingan DES dan AES}
\label{tab:des-aes}
\begin{tabular}{@{}lll@{}}
\toprule
\textbf{Aspek} & \textbf{DES} & \textbf{AES} \\
\midrule
Ukuran blok & 64 bit & 128 bit \\
Ukuran kunci & 56 bit & 128, 192, 256 bit \\
Jumlah putaran & 16 (Feistel) & 10, 12, 14 \\
Struktur & Feistel network & SPN (Substitution-Permutation) \\
Status & Tidak disarankan & Standar saat ini \\
\bottomrule
\end{tabular}
\end{table}

\subsection{DES}
Data Encryption Standard (DES) dikembangkan IBM dan diadopsi NIST pada 1977 \cite{nistfips46}. Menggunakan blok 64 bit dan kunci 56 bit. Struktur: 16 putaran Feistel, setiap putaran menggunakan S-box dan permutasi.

DES dianggap tidak aman untuk aplikasi modern karena ukuran kunci 56 bit terlalu kecil terhadap serangan brute force. Triple DES (3DES) menggunakan DES tiga kali dengan kunci berbeda untuk meningkatkan keamanan \cite{stallings2017,schneier1996}.

\subsection{AES}
Advanced Encryption Standard (AES) menggantikan DES pada 2001 \cite{nistfips197}. Menggunakan blok 128 bit dan kunci 128, 192, atau 256 bit. AES-128 adalah pilihan paling umum.

Struktur: matriks state 4$\times$4 byte. Setiap putaran meliputi SubBytes (S-box), ShiftRows, MixColumns, dan AddRoundKey. AES-128 memiliki 10 putaran. AES dirancang untuk resistan terhadap serangan diferensial dan linear \cite{nistfips197,stallings2017,schneier1996}.
